\documentclass[a4paper,12pt]{article}
\usepackage[utf8]{inputenc}
\usepackage[T1]{fontenc}
\usepackage[english]{babel}
\usepackage{geometry}
\geometry{left=2.5cm,right=2.5cm,top=2.5cm,bottom=2.5cm}

\usepackage{lmodern}
\usepackage{xcolor}
\definecolor{titleblue}{RGB}{0,51,140}

\usepackage{hyperref}

\pagestyle{empty}

\begin{document}

\begin{flushleft}
    \textbf{Ismail AISSAOUI} \\
    State Engineer in Artificial Intelligence \\
    ismail.aissaoui.pro@gmail.com \\
    +213 660 70 77 96 \\
    \href{https://ismail-aissaoui.vercel.app}{ismail-aissaoui.vercel.app}
\end{flushleft}

\begin{flushright}
    To the PhD Selection Committee \\
    \textbf{EDT Program - Université de Lille / INSA Rennes} \\
    Attn: Prof. Romain Rouvoy and Dr. Quentin Perez \\
    \medskip
    May 1, 2026
\end{flushright}

\bigskip

\noindent \textbf{Subject: Application for PhD: Improving Sustainability in the Physical-to-Digital Twin Continuum (Rank 8 - PC3)}

\bigskip

Dear Prof. Rouvoy and Dr. Perez,

As a State Engineer in Artificial Intelligence from ESI-SBA, I am writing to express my enthusiastic application for the PhD position in "Improving Sustainability in the Physical-to-Digital Twin Continuum" within the Engineering Digital Twin (EDT) national program. This application is motivated by my expertise in resource-constrained AI deployment and its relevance for your research on sustainable middleware.

Currently, I am completing my thesis at Sonatrach, where I have developed a **Physics-Informed Digital Twin** for gas turbine monitoring. A critical component of my work is the deployment of high-performance hybrid models (Mamba-SSM/xLSTM) onto resource-constrained industrial edge devices. This experience has confronted me with the direct trade-off between model precision and energy consumption. I have spent significant time optimizing the "energy footprint" of my surrogates to ensure they can provide real-time feedback without overwhelming the local compute infrastructure.

I am particularly interested in your project's focus on **adaptive deployment** and lifecycle sustainability. My technical background in **Distributed AI** and high-dimensional optimization aligns perfectly with the goal of creating a middleware reasoning framework that can dynamically select the right model fidelity and execution platform (Cloud vs. Edge) based on energy constraints. I am eager to apply my skills in software optimization and performance benchmarking to help create the **Artemis platform's** sustainable middleware layer, ensuring that digital twins are not only high-performing but also environmentally responsible.

My background in distributed AI and my experience with production-grade PyTorch and performance benchmarking provide me with the necessary tools to contribute to the **Artemis platform** and the broader EDT ecosystem. I am a highly autonomous researcher, comfortable bridging the gap between hardware constraints and sustainable software architecture.

Thank you for your time and consideration. I look forward to the possibility of discussing how my experience in adaptive deployment and energy-efficient AI can contribute to the success of your research group at the University of Lille and INSA Rennes.

Sincerely,

\bigskip

\textbf{Ismail AISSAOUI}

\end{document}
